\section{An explicit-density alternative: signed beta-prime weights}
\label{sec:beta-prime}
The learning proof uses only the scalar moment and comparison-law estimates, not infinite divisibility itself. A simple explicit density satisfies the same estimates.

Fix $p>2$ and $0<c_0\le1$. For fan-in $F$, put $a_F=c_0/F$ and give each coordinate the independent symmetric density
\begin{equation}
 \varrho_F(w)=
 \frac{|w|^{a_F-1}(1+|w|)^{-a_F-p}}
      {2\,\mathrm B(a_F,p)},\qquad w\ne0,
 \label{eq:beta-prime-prior}
\end{equation}
where $\mathrm B$ is Euler's beta function. Equivalently, $|W|$ has the beta-prime law with shape parameters $(a_F,p)$ and its sign is an independent fair coin. This is an absolutely continuous full-support distribution with no zero atom. Its density is sharply concentrated near zero and has a fixed polynomial tail. Fan-in changes its first shape parameter, not the architecture.

\begin{proposition}[The same learning guarantees for an explicit density]
\label{prop:beta-prime}
All circuit-neighborhood, expected-risk and Boolean online-prediction conclusions of this paper remain valid, with modified fixed constants, if every L\'evy increment coordinate is replaced by~\eqref{eq:beta-prime-prior}. The stable-teacher condition requires $A c_0/(p-1)\le1/2$ for unbounded activations, and the ReLU Boolean refinement requires $c_0/(p-1)\le1/2$.
\end{proposition}
\begin{proof}
We verify the scalar interface. Write $a=a_F\in(0,1]$ and
$\mathcal B=\mathrm B(a,p)$. The beta integral gives
\[
 \E|W|=\frac{a}{p-1},\qquad
 \E W^2=\frac{a(a+1)}{(p-1)(p-2)}.
\]
Thus the native first absolute and second moments are at most $C/F$.
Also $c_p/a\le\mathcal B\le1/a$, with fixed $c_p>0$. For the lower bound integrate the beta-prime density kernel on $(0,1)$; for the upper bound use the beta integral on $(0,1)$ with factor $(1-u)^{p-1}\le1$.

For $0<\epsilon\le1$, let $P_\epsilon=\Pr(|W|\le\epsilon)$.
The integral above $\epsilon$ gives
\[
 1-P_\epsilon\le Ca\{1+\log(1/\epsilon)\}.
\]
Indeed the integral on $(\epsilon,1)$ is at most
$\log(1/\epsilon)$, the integral on $(1,\infty)$ is at most $1/p$,
and $\mathcal B^{-1}\le Ca$.
A second bound, obtained by integrating below $\epsilon$, is
\[
 P_\epsilon\ge
 \frac{2^{-p-1}\epsilon^a}{a\mathcal B}
 \ge2^{-p-1}\epsilon^a.
\]
When $Ca\{1+\log(1/\epsilon)\}\le1/2$, the first bound and
$-\log(1-u)\le2u$ give the desired logarithmic conditioning cost.
Otherwise the second bound gives
$-\log P_\epsilon\le a\log(1/\epsilon)+(p+1)\log2$;
the fixed additive constant is absorbed by
$a\{1+\log(1/\epsilon)\}$ in this case. Therefore
\begin{equation}
 -\log P_\epsilon\le\frac C F\log(C/\epsilon).
 \label{eq:beta-prime-small-ball}
\end{equation}

Condition the coordinate directly on $|W|\le\epsilon$.
Symmetry preserves mean zero. For $r=1,2$, the ratio of its truncated moment integral to the lower bound on $P_\epsilon$ gives
\[
 \E[|W|^r\mid |W|\le\epsilon]
 \le 2^{p+1}\frac{a}{a+r}\epsilon^r
 \le C\epsilon^r/F.
\]
Together with~\eqref{eq:beta-prime-small-ball}, these are exactly the background support, entropy and moment estimates used in the network proof.

For a designated coefficient $u$, $|u|\le B$, use the unbiased separated interval law from Lemma~\ref{lem:levy-designated}, without adding a small-jump remainder. Its support lies in
$\{\epsilon<|w|\le B+8\}$, its density is at most $1/\epsilon$,
and its variance around $u$ is at most $C\epsilon^2$.
On this support,
\[
 \varrho_F(w)\ge
 \frac{a}{2(B+8)(B+9)^{p+1}}.
\]
Hence its relative entropy to $\varrho_F$ is at most
$\log F+C\log(C(B+8)/\epsilon)$.
All designated-support and unbiasedness properties are unchanged.

The remaining embedding, moment and statistical proofs use only these scalar estimates. They therefore apply verbatim. The native row's first absolute moment constant is $c_0/(p-1)$, giving the displayed intensity conditions. No Poisson representation is needed for this alternative.
\end{proof}

The density behaves near zero like $|w|^{-1+c_0/F}$, with integrable central concentration and no exact sparsity atoms. It can be viewed as continuous shrinkage whose strength is tied to fan-in. This is a mathematical prior construction, not an empirical claim about trained-network weights.

Its wide-row limit has the same L\'evy mechanism. For fixed $w\ne0$,
\[
 F\varrho_F(w)\longrightarrow
 \frac{c_0}{2|w|(1+|w|)^p}.
\]
Dominated convergence against $\cos(tw)-1$, using the quadratic bound near zero and the integrable tail bound, shows that a sum of $F$ such independent coordinates converges to the non-Gaussian infinitely divisible law with this L\'evy measure. Its variance remains of order one. Exact finite-width infinite divisibility is special to increment laws and is not needed for any learning conclusion.
